\section{Preserved Identification and Native-Task Protocol}
\label{app:prospective}
This appendix documents the earlier identification and native-task protocol, including branches that stopped or never reached confirmation. The completed follow-up analyses are reported in \cref{sec:pending-studies}. The earlier calibration and native branches documented here retain their own recorded execution states. Structural recovery, excluded development analyses and later descriptive outcomes have different evidential roles.

\paragraph{Completed characterization and computational reproduction.}
The revised BF16/F64 synthetic calibration completed, while the original calibration failure remains recorded. The excluded pilot completed structural eligibility and power calculations; structural power remained limited and native preparation stopped at clean competence. The separate coordinate characterization was specified after excluded development work. Its results were subsequently inspected and manually summarized descriptively. Reproducing those computations does not restore an untouched test set or establish independent confirmation.

The subsequent computational reproduction completed the fixed compact/context panel, rank-aware historical continuation, peer adapters, factorial comparisons, word and structural cross-model transport, product actions and constructive interfaces (\cref{app:reproduction}). These branches have distinct inputs and evidential roles. Their completion neither overrides the native competence stop nor authorizes the blocked confirmation branches. The original characterization remains identifiable: it did not itself include these later extensions. Original output flags, failed units, numerical tables and source identities are preserved; reviewed derivative tables have separate provenance.

The characterization uses the six pinned base models, four readouts and training-only $q=16$ PCA. Its family/crossed 10,000-replicate percentile intervals (seed 314159) describe sensitivity of a fixed panel; they are not new multiplicity-adjusted hypothesis tests or model-population confidence statements. The completed compact/context extension reports descriptive fixed-panel means without new significance decisions. Verification covers archived integrity, stored selections, historical numerical agreement and representative analytic refits. Raw pretrained forwards, raw-to-PCA fits and complete neural retraining are outside that local verification.

\subsection{Earlier protocol as specified, including unexecuted branches}
The following paragraphs describe the preserved design, not a list of completed experiments or instructions for the present analysis. Execution status is reported in \cref{tab:reproduction-status}. In particular, the proposed native interventions and confirmatory decision rules below did not acquire confirmatory outcomes.

\paragraph{Sampling and partitions.}
The default structural protocol specified 48 training, 16 validation and 32 test scenario families, with two paired paraphrases per family and 12 additional pilot families. Its declared compositional generator sampled scenarios under a fixed collection of grammar roots; this supports inference conditional on that generator, not a claim of 96 independently authored grammars. The design allowed external authored banks to be imported with family and grammar-root identities. The protocol called for sealing a maximum candidate test pool before pilot analysis. The 32-family default was the smallest allowable prefix. Prospective power was to select 32, 48 or 64 families before collecting the corresponding activations. The maximum pool contained 24 new entity triples, with crossed power intended to select 16 or 24. All states of a scenario remain together. The proposed per-root and leave-one-root-out summaries would assess sensitivity to the fixed grammar generator, not transfer to independently authored grammars; their specification alone is not a completed result. The six historical checkpoints and four relative depths are fixed conditions rather than independent model-population draws.

\paragraph{New entity transfer.}
Alongside $C_7$ and the historical four-corner $D_4$ task, an $S_3$ role-binding task crosses scenario families with 48 disjoint entity triples, split 24/8/16. Each triple has six assignments. Training uses training-family by training-triple cells; validation uses validation by validation. The principal transfer block contains new families and new triples. A single fixed final-token readout measures the complete clause; concatenating role-aligned hidden vectors would engineer the desired action and is excluded from the primary measurement. Zero-shot evaluation uses no target-triple activation centroids. A separately charged equal-support arm provides four disjoint support contexts per assignment state to all applicable methods.

\paragraph{Strong alternative explanations.}
Training centroids $\mu_i$ define $z-\mu_i+\mu_{g(i)}$ and scalar residual-retention baselines. Their state indices update under composition. Oracle-source and construction-decoded-source conditions are separate. Independent linear ridge is the primary prediction procedure; ridge on training context-difference pairs is the co-primary transport procedure. Affine, Procrustes, residual, Joint and validation-selected local-displacement variants are secondary. The local field is evaluated against actual held-out targets and reported separately from the analytic linear-estimator null calibration. Regularization is selected on validation, with penalty $\rho=\lambda\operatorname{tr}(X^\top X)/q$ and a declared grid including a pseudoinverse candidate. The $S_3$ comparison additionally fits a compositional potential field to construction entity embeddings and role metadata, without held-out target activations.

For paraphrases $t,t'$ within family $f$, the protocol specified transport of $u_i=z_{t,i}-z_{t',i}$ to $v_{g,i}=z_{t,g(i)}-z_{t',g(i)}$, against zero and validation-selected scalar identity transport. Target and difference losses use common training-only energy scales. Low training context energy makes the difference endpoint uninformative and is retained as such. Exact coherent fixed-feature action reassignments provide 120 cyclic image families and three four-corner dihedral image families. For the latter, the smallest tail fraction is $1/3$ because the count includes the registered image itself. This earlier reference refitted predictive maps for each coherent image on fixed features and scored complete state/action predictions, including identity and fixed-point pairs. Its prediction-loss ranks are distinct from the later spectral-budget rankings; neither is automatically a semantic randomization $p$-value.

\paragraph{Calibration and uncertainty.}
The protocol specified six synthetic regimes to test category-only realizations, shared context, independent state--context noise, genuinely equivariant interaction, family-specific transport and unsupported ambient extensions. Calibration was intended to assess the actual fitted-estimator decision rule, while keeping statistical error distinct from semantic nonidentification. The proposed primary $C_7/D_4$ inference would aggregate fixed models and depths within each family and use family-level tests with Holm correction. Its eight superiority/equivalence hypotheses were allocated familywise level .04; the four zero-shot $S_3$ hypotheses were allocated .01 with a separately calibrated crossed family/orbit analysis. These allocations are preserved design choices, not achieved confirmatory error guarantees for the descriptive results. The proposed normalized superiority margin is .01 and equivalence half-width .005. These are research choices, not conference thresholds. Bootstrap results are sensitivity analyses; overlapping split seeds do not multiply the sample size.

\paragraph{Native computation.}
The proposed causal study started from Llama-3.1-8B and Qwen2.5-7B, subject to a development-only clean-competence gate. Its design used genuine weekday programs and a separate free-eight-state $D_4$ task on an asymmetrically marked square, rather than the four-corner state action of the representation study. For prefix $p$, generator edit $g$ and native continuation $v$, the counterfactual answer is $v(g(p(i)))$. A natural counterpart starts from $p^{-1}(g(p(i)))$, so applying the prefix reaches the same intermediate state $g(p(i))$. The model executes $v$. Externally composing an edit is not described as the model executing those matrices. A candidate intermediate site must carry the proposed variable without inserting its ground-truth answer into the prompt.

For span means $\bar h,\bar h^*$, the proposed primary hook was $h_j'=h_j+\alpha(\bar h^*-\bar h)$, which preserves token deviations from the span mean. Preserving a declared complement additionally requires the mean displacement to lie within the chosen intervention subspace. This proposed hook differs from the historical mean-overwrite intervention in \cref{app:behavior}. The planned controls included zero-delta and historical mean-overwrite shams, natural donors, source-conditioned means, shrinkage, affine/Procrustes/local fields and fixed-seed random carriers. Earlier consumer sites were designated primary and late patches positive controls. The plan selected the best trainable baseline on validation and specified complete-answer, sequence-margin, invalid-output and collateral-effect scores. The proposed primary accuracy contrast used held-out continuation words of length six to eight, with shorter words assigned a separate sensitivity condition. The excluded pilot uses the test panel's long/short allocation. A five-point observed accuracy advantage with an adjusted positive bound is a proposed practical target; it does not establish a unique native algorithm. Operation-only, joint-variable and pathway interventions require their separately defined carriers and tests. A proposed secondary capacity panel would fit ranks two, four and eight at the frozen primary sites, freezing its construction fits and validation-selected baselines before test outcomes. Its intended rank-two action-map comparison with a rank-two calendar plane would not replace the primary rank-sixteen design. The blocked native confirmation provides no outcome for this proposed comparison.

\paragraph{Prospective peer comparison (unexecuted).}
The protocol proposed comparing calendar-plane interventions \citep{engels2025notall}, conditional-mean relation steering \citep{arora2026causal}, and manifold steering \citep{wurgaft2026manifold} on the same sealed cases. That comparison did not run. \Cref{app:reproduction} reports separate method-inspired Engels and Wurgaft anchors and an Arora-inspired factorial study, each with its own tasks and scope.

The proposed comparison would match intervention masks, treatment of the complement, rank, supplied labels and tuning budgets. The published calendar-plane intervention mean-ablates the complement \citep[Section~5.1]{engels2025notall}. The reported manifold intervention preserves the PCA complement, while its straight-line comparator replaces the full activation \citep[Appendix~A.6]{wurgaft2026manifold}. A comparison that uses one complement rule would therefore adapt at least one of the original protocols. The conditional-mean decomposition would use a factorial ensemble with all factor combinations \citep[Section~2]{arora2026causal}. The completed cross-model word comparison in \cref{app:reproduction} uses a distinct information condition: Hu's method predicts target-model displacement using the observed source-model displacement for the test word \citep[Equation~12]{hu2026language}.

A separate optional comparison could reassign symbols to group elements in each prompt, following the task construction of \citet[Section~2]{todd2026algebra}. It would ask whether predictive recovery survives when tokens no longer have fixed group meanings. This comparison was outside the executed studies.

\paragraph{Reporting policy for the prospective protocol.}
The protocol called for retaining negative, infeasible and inconclusive conditions. Under its interpretation rules, a category or local baseline accounting for the effect would narrow the structural claim. Success restricted to late patches would support a control interface. A claim about native use would require counterfactual fidelity and specificity within the tested abstraction. A third training lineage was designated as a separate prospective replication, with the primary panel kept fixed. The later Mistral fitted-bank anchor is not that prespecified independent replication. \Cref{tab:reproduction-status} lists the recorded outcomes, including the native clean-competence stop. The blocked confirmation branches have no measured outcomes.

